% ============================================
% MEM-09
% Gerado de molambudos.md — build_miolo.py
% ============================================
\fragtitle{mem}{M}{MEM}{09}{Os Primeiros Anos}
\markboth{MEM-09 --- Os Primeiros Anos}{}

\noindent \textit{O menino que não falava}

\noindent \textbf{1920-1930 --- O menino que não falava}

\vspace{0.5em}\noindent\rule{\textwidth}{0.5pt}\vspace{0.5em}

\noindent Nos primeiros anos no Colônia, eu não falava.

\noindent Não porque não pudesse. Porque não valia a pena.

\noindent Falar significava responder. Responder significava explicar. Explicar significava repetir a história. E a história --- a vala, a criatura, o olho amarelo --- ninguém ia entender.

\noindent Então eu ficava quieto.

\noindent Os enfermeiros gostavam de paciente quieto. Paciente quieto não dá trabalho. Paciente quieto não toma tempo. Paciente quieto é como móvel --- está ali, ocupa espaço, mas não exige nada.

\noindent Eu era um móvel.

\noindent Sentado no canto do pátio. Olhando para o chão. O dia inteiro. Todos os dias.

\noindent A criatura dormia dentro de mim. De vez em quando ela acordava, se espreguiçava, e eu sentia uma pontada no estômago. Ela estava crescendo.

\noindent \textit{\textquotedblleft{}Tu precisa comer mais\textquotedblright{}}, ela disse.

\noindent \textit{\textquotedblleft{}Não tem comida.\textquotedblright{}}

\noindent \textit{\textquotedblleft{}Tem. Tu só não está vendo.\textquotedblright{}}

\noindent \textit{\textquotedblleft{}Onde?\textquotedblright{}}

\noindent \textit{\textquotedblleft{}Nos outros.\textquotedblright{}}

\noindent Eu não respondi.

\noindent Mas eu entendi.

\vspace{0.5em}\noindent\rule{\textwidth}{0.5pt}\vspace{0.5em}

\noindent Em 1922, eu matei o Wagner.

\noindent Depois disso, os enfermeiros começaram a me tratar diferente. Não chegavam perto. Deixavam a comida na porta da cela. Evitavam olhar nos meus olhos.

\noindent O isolamento virou minha casa.

\noindent Uma cela de três metros por dois. Uma cama de ferro. Um balde. Uma janela gradeada que dava para o corredor. Eu ficava ali, sentado, escrevendo no diário.

\noindent As horas passavam.

\noindent Os dias passavam.

\noindent Os anos passavam.

\noindent Eu perdia a conta do tempo. O diário virava meu relógio. Cada página virada era um dia. Cada página escrita era uma noite.

\noindent A criatura falava cada vez mais.

\noindent \textit{\textquotedblleft{}Escreve.\textquotedblright{}}

\noindent \textit{\textquotedblleft{}Sobre o quê?\textquotedblright{}}

\noindent \textit{\textquotedblleft{}Sobre tudo. Sobre a vala. Sobre o curral. Sobre a fome. Escreve para que alguém leia.\textquotedblright{}}

\noindent \textit{\textquotedblleft{}Quem vai ler?\textquotedblright{}}

\noindent \textit{\textquotedblleft{}Alguém. Um dia. Alguém vai abrir este diário e vai ver.\textquotedblright{}}

\noindent \textit{\textquotedblleft{}Ver o quê?\textquotedblright{}}

\noindent \textit{\textquotedblleft{}A verdade.\textquotedblright{}}

\noindent \textit{\textquotedblleft{}Qual verdade?\textquotedblright{}}

\noindent \textit{\textquotedblleft{}A verdade que ninguém quer ver. Que o mundo é uma vala. Que a gente nasce, passa fome, morre, e é esquecido. Mas o diário não esquece. O diário guarda.\textquotedblright{}}

\noindent E eu escrevia.

\noindent Página após página.

\noindent Ano após ano.

\noindent Escrevia até os dedos sangrarem.

\noindent Escrevia até a tinta acabar --- e aí escrevia com sangue.

\noindent O diário crescia comigo.

\noindent Ele sabia coisas que eu nem sabia que sabia.

\noindent Às vezes eu abria numa página antiga e não reconhecia a letra. Achava que era de outra pessoa. Mas era a minha. Mão minha. Letra minha.

\noindent Mas as palavras...

\noindent As palavras não eram minhas.

\noindent Eram dela.

\vspace{0.5em}
\noindent\textit{\textbf{↪ Rotas:} \rota{MEM-08} \rotasep \rota{MEM-10}}